\documentclass[11pt,a4paper]{article}
\usepackage[T1]{fontenc}
\usepackage[margin=20mm]{geometry}
\usepackage{xcolor}
\usepackage{listings}
\usepackage[hidelinks]{hyperref}
\pagestyle{plain}
\definecolor{codebg}{RGB}{246,248,250}
\lstset{basicstyle=\ttfamily\small,backgroundcolor=\color{codebg},
  frame=single,framerule=0.3pt,breaklines=true,columns=fullflexible,
  xleftmargin=4pt,xrightmargin=4pt,aboveskip=7pt,belowskip=7pt}
\begin{document}
\begin{center}
{\Large\bfseries Required One-Time FAST-LIO2 Installation}\par
\vspace{3pt}
{\small Husky GNSS-Denied Navigation --- Team Instructions}
\end{center}

\section*{Before running the Bayland command}
FAST-LIO2 build products are excluded from Git. Cloning or pulling the
repository supplies the source but does not create the executable. Perform this
step on every new checkout or machine, from a fresh terminal.

\begin{lstlisting}
cd <Journal_Work>/HIL/verified_stack
source /opt/ros/jazzy/setup.bash
bash hardware/build_fastlio.sh
source FASTLIO2Experiment/install/setup.bash
\end{lstlisting}

The build must finish without error. Do not start the experiment when
\path{FASTLIO2Experiment/install/setup.bash} is missing.

\section*{Mandatory verification}
\begin{lstlisting}
test -x FASTLIO2Experiment/install/fast_lio/lib/fast_lio/fastlio_mapping \
  && echo "Correct FAST-LIO2 build exists"

ros2 pkg prefix fast_lio
\end{lstlisting}

The package prefix must begin with:

\begin{lstlisting}
<Journal_Work>/HIL/verified_stack/FASTLIO2Experiment/install
\end{lstlisting}

If it resolves to \texttt{HIL/verified\_stack/install}, another workspace, or
an older installation, stop. Open a fresh terminal, source ROS Jazzy, rebuild
with the command above, and repeat the verification. The same navigation
command does not reproduce the same experiment when ROS selects a different
FAST-LIO2 executable.

\section*{Then run the verified profile}
Continue with Section 1 only after the local executable and package prefix have
both been verified. Keep the complete result folder from every run. In
\texttt{fast\_lio2.log}, the launched executable must also be under
\texttt{HIL/verified\_stack/FASTLIO2Experiment/install/fast\_lio/}.

\vfill
\noindent\textbf{Fail-safe rule:} a missing local FAST-LIO2 build is a setup
failure, not permission to fall back to an executable from another ROS overlay.
\end{document}
